% !TEX root = ../main.tex
%% Corrigés — Chapitre 11 — Analyse asymptotique
%% Résultats vérifiés par verifications/chapitre_11.py.

\corrigesChapitre{11}{Analyse asymptotique}
\corrigesSouples

\begin{blocLivre}{cor-11-1}
\begin{corrige}{1}{Premiers équivalents}
\demarche{Tout repose sur un quotient dont on connaît la limite.}
\begin{questions}
  \item $\dfrac{\sin x}{x}\to1$ (taux d'accroissement de $\sin$ en $0$) ;
    $\dfrac{\tan x}{x}=\dfrac{\sin x}{x}\cdot\dfrac{1}{\cos x}\to1$ ;
    enfin $1-\cos x=2\sin^{2}\dfrac{x}{2}\sim2\Bigl(\dfrac{x}{2}\Bigr)^{2}
    =\dfrac{x^{2}}{2}$.
  \item $\dfrac{\ln(1+x)}{x}\to\ln'(1)=1$ et
    $\dfrac{\mathrm{e}^{x}-1}{x}\to\exp'(0)=1$ : deux taux d'accroissement.
  \item $\dfrac{3x^{3}+x^{2}}{3x^{3}}=1+\dfrac{1}{3x}\to1$ en $\pm\infty$ ;
    $\dfrac{3x^{3}+x^{2}}{x^{2}}=1+3x\to1$ en $0$ : terme dominant d'un
    côté, terme de plus bas degré de l'autre.
  \item $\dfrac{\sqrt{x^{2}+x}}{x}=\sqrt{1+\dfrac1x}\to1$ en $+\infty$.
    En $-\infty$, $\sqrt{x^{2}+x}=\lvert x\rvert\sqrt{1+\tfrac1x}$ donc
    $\sqrt{x^{2}+x}\sim\res{-x}$ (le signe compte !).
\end{questions}
\end{corrige}
\end{blocLivre}

\begin{blocLivre}{cor-11-2}
\begin{corrige}{2}{Limites par équivalents}
\demarche{Remplacer chaque facteur par un équivalent — jamais un terme d'une somme.}
\begin{questions}
  \item $\dfrac{\sin3x}{\tan5x}\sim\dfrac{3x}{5x}$ :
    limite $\res{\dfrac35}$.
  \item $1-\cos x\sim\dfrac{x^{2}}{2}$ : limite $\res{\dfrac12}$.
  \item $\ln\Bigl(1+\dfrac1x\Bigr)\sim\dfrac1x$ (composition à droite par
    $\tfrac1x\to0$) : limite $\res{1}$.
  \item $\mathrm{e}^{2x}-1\sim2x$ et $\sin x\sim x$ :
    limite $\res{2}$.
\end{questions}
\end{corrige}
\end{blocLivre}

\begin{blocLivre}{cor-11-3}
\begin{corrige}{3}{DL usuels : échauffement}
\demarche{Les formules du cours, tronquées à l'ordre demandé.}
\begin{questions}
  \item $\mathrm{e}^{x}=1+x+\dfrac{x^{2}}{2}+\dfrac{x^{3}}{6}+o\bigl(x^{3}\bigr)$
  \item $\cos x=1-\dfrac{x^{2}}{2}+\dfrac{x^{4}}{24}+o\bigl(x^{4}\bigr)$
  \item $\ln(1+x)=x-\dfrac{x^{2}}{2}+\dfrac{x^{3}}{3}+o\bigl(x^{3}\bigr)$
  \item $\sqrt{1+x}=1+\dfrac{x}{2}-\dfrac{x^{2}}{8}+o\bigl(x^{2}\bigr)$
    (cas $\alpha=\tfrac12$)
  \item $\dfrac{1}{1-x}=1+x+x^{2}+x^{3}+o\bigl(x^{3}\bigr)$
  \item $\dfrac{1}{1+x}=1-x+x^{2}-x^{3}+o\bigl(x^{3}\bigr)$
    (remplacer $x$ par $-x$)
\end{questions}
\end{corrige}
\end{blocLivre}

\begin{blocLivre}{cor-11-4}
\begin{corrige}{4}{Produits de développements limités}
\demarche{Multiplier les parties régulières et tronquer sans états d'âme.}
\begin{questions}
  \item $\Bigl(1+x+\dfrac{x^{2}}{2}+\dfrac{x^{3}}{6}\Bigr)
    \Bigl(x-\dfrac{x^{3}}{6}\Bigr)
    =x+x^{2}+\Bigl(\dfrac12-\dfrac16\Bigr)x^{3}+o\bigl(x^{3}\bigr)$, d'où
    \[
    \mathrm{e}^{x}\sin x=\res{x+x^{2}+\frac{x^{3}}{3}+o\bigl(x^{3}\bigr)}.
    \]
  \item Par produit :
    $\Bigl(1-\dfrac{x^{2}}{2}+\dfrac{x^{4}}{24}\Bigr)^{2}
    =1-x^{2}+\Bigl(\dfrac14+\dfrac{1}{12}\Bigr)x^{4}+o\bigl(x^{4}\bigr)
    =1-x^{2}+\dfrac{x^{4}}{3}+o\bigl(x^{4}\bigr)$.
    Par linéarisation : $\cos2x=1-2x^{2}+\dfrac{2x^{4}}{3}+o\bigl(x^{4}\bigr)$
    (remplacer $x$ par $2x$), et
    $\dfrac{1+\cos2x}{2}=\res{1-x^{2}+\frac{x^{4}}{3}+o\bigl(x^{4}\bigr)}$ :
    les deux calculs concordent.
\end{questions}
\end{corrige}
\end{blocLivre}

\begin{blocLivre}{cor-11-5}
\begin{corrige}{5}{Quotients et composées}
\demarche{Pour diviser : factoriser la valeur en 0, puis composer avec
$\frac{1}{1+u}=1-u+u^{2}-\dots$}
\begin{questions}
  \item $\dfrac{1}{\cos x}=\dfrac{1}{1-\frac{x^{2}}{2}+o\left(x^{3}\right)}
    =1+\dfrac{x^{2}}{2}+o\bigl(x^{3}\bigr)$, puis
    $\tan x=\Bigl(x-\dfrac{x^{3}}{6}\Bigr)\Bigl(1+\dfrac{x^{2}}{2}\Bigr)
    +o\bigl(x^{3}\bigr)=\res{x+\frac{x^{3}}{3}+o\bigl(x^{3}\bigr)}$.
  \item $1+\mathrm{e}^{x}=2\Bigl(1+\dfrac{x}{2}+\dfrac{x^{2}}{4}+o\bigl(x^{2}\bigr)\Bigr)$,
    donc avec $u=\dfrac{x}{2}+\dfrac{x^{2}}{4}$ :
    $\dfrac{1}{1+\mathrm{e}^{x}}=\dfrac12\bigl(1-u+u^{2}\bigr)+o\bigl(x^{2}\bigr)
    =\dfrac12\Bigl(1-\dfrac{x}{2}-\dfrac{x^{2}}{4}+\dfrac{x^{2}}{4}\Bigr)+o\bigl(x^{2}\bigr)
    =\res{\frac12-\frac{x}{4}+o\bigl(x^{2}\bigr)}$ :
    le coefficient de $x^{2}$ est nul (la fonction $x\mapsto\frac{1}{1+\mathrm{e}^{x}}-\frac12$
    est impaire).
  \item $\ln(\cos x)=\ln(1+u)$ avec $u=-\dfrac{x^{2}}{2}+o\bigl(x^{3}\bigr)$ :
    comme $u^{2}=O\bigl(x^{4}\bigr)$,
    $\ln(\cos x)=u+o\bigl(x^{3}\bigr)=\res{-\frac{x^{2}}{2}+o\bigl(x^{3}\bigr)}$.
\end{questions}
\end{corrige}
\end{blocLivre}

\begin{blocLivre}{cor-11-6}
\begin{corrige}{6}{Primitiver un développement limité}
\demarche{On primitive terme à terme, sans oublier la valeur en 0.}
\begin{questions}
  \item $\dfrac{1}{1+x^{2}}=1-x^{2}+x^{4}+o\bigl(x^{4}\bigr)$, et
    $\arctan0=0$ :
    $\arctan x=\res{x-\frac{x^{3}}{3}+\frac{x^{5}}{5}+o\bigl(x^{5}\bigr)}$.
  \item $\dfrac{1}{1+x}=1-x+x^{2}-x^{3}+o\bigl(x^{3}\bigr)$ et $\ln1=0$ :
    $\ln(1+x)=x-\dfrac{x^{2}}{2}+\dfrac{x^{3}}{3}-\dfrac{x^{4}}{4}+o\bigl(x^{4}\bigr)$.
  \item Composition licite car $x^{2}\to0$ :
    $\ln\bigl(1+x^{2}\bigr)=x^{2}-\dfrac{x^{4}}{2}+o\bigl(x^{4}\bigr)$ —
    directement, en substituant $x^{2}$ dans $\ln(1+u)=u-\frac{u^{2}}{2}+o\bigl(u^{2}\bigr)$.
\end{questions}
\end{corrige}
\end{blocLivre}

\begin{blocLivre}{cor-11-7}
\begin{corrige}{7}{Courbe et tangente}
\demarche{La tangente, c'est le DL à l'ordre 1 ; la position, le premier terme suivant.}
\begin{questions}
  \item $\mathrm{e}^{x}=1+x+\dfrac{x^{2}}{2}+o\bigl(x^{2}\bigr)$ : tangente
    $y=1+x$, et $\dfrac{x^{2}}{2}>0$ (ordre pair, coefficient positif) :
    courbe \res{\text{au-dessus}} au voisinage de $0$.
  \item $\ln(1+x)=x-\dfrac{x^{2}}{2}+o\bigl(x^{2}\bigr)$ : tangente $y=x$,
    courbe \res{\text{au-dessous}}.
  \item $\sin x=x-\dfrac{x^{3}}{6}+o\bigl(x^{3}\bigr)$ : tangente $y=x$,
    premier terme correctif d'ordre \emph{impair} : la courbe
    \res{\text{traverse}} sa tangente (point d'inflexion).
\end{questions}
\etape{Global pour f}
Soit $\varphi(x)=\mathrm{e}^{x}-1-x$ : $\varphi'(x)=\mathrm{e}^{x}-1$ est
négative sur $\R_{-}$, positive sur $\R_{+}$, donc $\varphi$ atteint son
minimum en $0$, où $\varphi(0)=0$ : $\mathrm{e}^{x}\geqslant1+x$ pour tout
$x\in\R$ (c'est aussi la convexité, chapitre 6).
\end{corrige}
\end{blocLivre}

\begin{blocLivre}{cor-11-8}
\begin{corrige}{8}{Équivalents de suites}
\demarche{Quantité conjuguée, composition à droite, et DL pour les sommes.}
\begin{questions}
  \item $\sqrt{n+1}-\sqrt{n}=\dfrac{1}{\sqrt{n+1}+\sqrt{n}}$, et
    $\sqrt{n+1}+\sqrt{n}\sim2\sqrt{n}$ :
    d'où $\sqrt{n+1}-\sqrt{n}\sim\res{\dfrac{1}{2\sqrt{n}}}$.
  \item $\sin\dfrac1n\sim\dfrac1n$ (car $\tfrac1n\to0$) ;
    $n^{2}+3n\sim n^{2}$ ;
    $\ln\Bigl(1+\dfrac{1}{n^{2}}\Bigr)\sim\dfrac{1}{n^{2}}$.
  \item $\ln u_n=n\ln\Bigl(1+\dfrac1n\Bigr)
    =n\Bigl(\dfrac1n+o\Bigl(\dfrac1n\Bigr)\Bigr)=1+o(1)\to1$, et par
    continuité de l'exponentielle, $u_n\to\res{\mathrm{e}}$. (On compose ici
    une \emph{limite}, ce qui est toujours permis.)
\end{questions}
\end{corrige}
\end{blocLivre}

\begin{blocLivre}{cor-11-9}
\begin{corrige}{9}{Limites plus délicates}
\demarche{Forme exponentielle pour les puissances, DL dès qu'une somme apparaît.}
\begin{questions}
  \item $(\cos x)^{1/x^{2}}=\exp\Bigl(\dfrac{\ln\cos x}{x^{2}}\Bigr)$ et
    $\ln\cos x\sim-\dfrac{x^{2}}{2}$ : l'exposant tend vers $-\dfrac12$,
    limite $\res{\mathrm{e}^{-1/2}}$.
  \item $(1+x)^{1/x}=\exp\Bigl(\dfrac{\ln(1+x)}{x}\Bigr)
    =\exp\Bigl(1-\dfrac{x}{2}+o(x)\Bigr)
    =\mathrm{e}\Bigl(1-\dfrac{x}{2}+o(x)\Bigr)$,
    donc $\dfrac{(1+x)^{1/x}-\mathrm{e}}{x}\to\res{-\dfrac{\mathrm{e}}{2}}$.
  \item $\tan x-x\sim\dfrac{x^{3}}{3}$ et $x-\sin x\sim\dfrac{x^{3}}{6}$
    (DL du cours) : limite $\res{2}$.
  \item Avec $u=\dfrac1x\to0$ :
    $x^{2}\ln\Bigl(1+\dfrac1x\Bigr)=\dfrac{u-\frac{u^{2}}{2}+o\left(u^{2}\right)}{u^{2}}
    =x-\dfrac12+o(1)$, donc
    $x-x^{2}\ln\Bigl(1+\dfrac1x\Bigr)\to\res{\dfrac12}$.
\end{questions}
\end{corrige}
\end{blocLivre}

\begin{blocLivre}{cor-11-10}
\begin{corrige}{10}{Développements en un point quelconque}
\demarche{Poser $x=a+h$ et tout ramener en $h=0$.}
\begin{questions}
  \item $\ln(1+h)=h-\dfrac{h^{2}}{2}+o\bigl(h^{2}\bigr)$ ;
    $\mathrm{e}^{1+h}=\mathrm{e}\,\mathrm{e}^{h}
    =\mathrm{e}+\mathrm{e}h+\dfrac{\mathrm{e}}{2}h^{2}+\dfrac{\mathrm{e}}{6}h^{3}+o\bigl(h^{3}\bigr)$.
  \item $\sqrt{4+h}=2\sqrt{1+\dfrac{h}{4}}
    =2\Bigl(1+\dfrac{h}{8}-\dfrac{h^{2}}{128}+o\bigl(h^{2}\bigr)\Bigr)
    =\res{2+\frac{h}{4}-\frac{h^{2}}{64}+o\bigl(h^{2}\bigr)}$.
  \item $\sin\Bigl(\dfrac{\pi}{6}+h\Bigr)
    =\dfrac12\cos h+\dfrac{\sqrt3}{2}\sin h
    =\res{\frac12+\frac{\sqrt3}{2}h-\frac{h^{2}}{4}+o\bigl(h^{2}\bigr)}$.
\end{questions}
\end{corrige}
\end{blocLivre}

\begin{blocLivre}{cor-11-11}
\begin{corrige}{11}{Prolongement et dérivabilité par DL}
\demarche{Un DL d'ordre 1 du prolongement, lu sur le DL d'ordre 2 du numérateur.}
\begin{questions}
  \item $\mathrm{e}^{x}-1=x+\dfrac{x^{2}}{2}+\dfrac{x^{3}}{6}+o\bigl(x^{3}\bigr)$,
    donc $f(x)=1+\dfrac{x}{2}+\dfrac{x^{2}}{6}+o\bigl(x^{2}\bigr)$ :
    $f\to1$, on prolonge par $f(0)=1$. Alors
    $\dfrac{f(x)-1}{x}=\dfrac12+\dfrac{x}{6}+o(x)\to\dfrac12$ :
    le prolongement est dérivable en $0$ et $\res{f'(0)=\dfrac12}$.
  \item $g(x)=\dfrac{\sin x}{x}=1-\dfrac{x^{2}}{6}+o\bigl(x^{2}\bigr)$ :
    prolongement $g(0)=1$, et $\dfrac{g(x)-1}{x}=-\dfrac{x}{6}+o(x)\to0$,
    donc $g$ est dérivable en $0$ avec $\res{g'(0)=0}$. Enfin
    $g(x)-1=-\dfrac{x^{2}}{6}+o\bigl(x^{2}\bigr)<0$ au voisinage épointé
    de $0$ : le point $0$ est un \res{\text{maximum local}} (strict).
\end{questions}
\end{corrige}
\end{blocLivre}

\begin{blocLivre}{cor-11-12}
\begin{corrige}{12}{Extremums : l'ordre 2… et au-delà}
\demarche{Le premier terme non nul du DL décide : pair et positif, pair et
négatif, ou impair.}
\begin{questions}
  \item $f(x)=x^{2}+x^{3}=x^{2}(1+x)$ : premier terme $x^{2}$, d'ordre pair
    et de coefficient $1>0$ : \res{\text{minimum local}}.
  \item $g(x)=x^{3}+x^{4}=x^{3}(1+x)$ : premier terme d'ordre \emph{impair},
    $g$ change de signe en $0$ ($g(-0{,}01)<0<g(0{,}01)$) :
    \res{\text{ni maximum ni minimum}}.
  \item $h'(0)=0$ et $h''(0)=-\cos0+1=0$ : l'ordre 2 ne décide pas. Mais
    \[
    h(x)=\cos x+\frac{x^{2}}{2}
    =1-\frac{x^{2}}{2}+\frac{x^{4}}{24}+\frac{x^{2}}{2}+o\bigl(x^{4}\bigr)
    =1+\frac{x^{4}}{24}+o\bigl(x^{4}\bigr),
    \]
    et $\dfrac{x^{4}}{24}>0$ : $0$ est un \res{\text{minimum local strict}}.
    La condition suffisante d'ordre 2 peut donc échouer ; le DL, lui,
    tranche toujours dès qu'un terme non nul apparaît.
\end{questions}
\end{corrige}
\end{blocLivre}

\begin{blocLivre}{cor-11-13}
\begin{corrige}{13}{Asymptotes par développement}
\demarche{Factoriser le terme dominant, poser $u=1/x$, développer à l'ordre
$u^{2}$ avant de remultiplier par $x$.}
\begin{questions}
  \item Pour $x>0$ :
    $f(x)=x\sqrt{1+u+u^{2}}$ avec $u=\dfrac1x$, et
    \[
    \sqrt{1+u+u^{2}}=1+\frac{u+u^{2}}{2}-\frac{(u+u^{2})^{2}}{8}+o\bigl(u^{2}\bigr)
    =1+\frac{u}{2}+\frac{3u^{2}}{8}+o\bigl(u^{2}\bigr),
    \]
    d'où $f(x)=x+\dfrac12+\dfrac{3}{8x}+o\Bigl(\dfrac1x\Bigr)$. La droite
    $\res{y=x+\dfrac12}$ est asymptote, et $\dfrac{3}{8x}>0$ : courbe
    au-dessus en $+\infty$.
  \item Pour $x<0$, $\sqrt{x^{2}}=-x$ :
    $f(x)=-x-\dfrac12-\dfrac{3}{8x}+o\Bigl(\dfrac1x\Bigr)$, asymptote
    $\res{y=-x-\dfrac12}$, et $-\dfrac{3}{8x}>0$ : la courbe est encore
    au-dessus. (Contrôle : $f(-10)=\sqrt{91}\approx9{,}539>9{,}5$.)
\end{questions}
\end{corrige}
\end{blocLivre}

\begin{blocLivre}{cor-11-14}
\begin{corrige}{14}{La formule de Stirling en action}
\demarche{Stirling est un équivalent : on l'injecte dans des quotients, et
son logarithme donne un développement.}
\begin{questions}
  \item $\displaystyle\binom{2n}{n}=\frac{(2n)!}{(n!)^{2}}
    \sim\frac{\sqrt{4\pi n}\,\bigl(\frac{2n}{\mathrm{e}}\bigr)^{2n}}
    {2\pi n\,\bigl(\frac{n}{\mathrm{e}}\bigr)^{2n}}
    =\frac{2^{2n}\sqrt{4\pi n}}{2\pi n}
    =\res{\frac{4^{n}}{\sqrt{\pi n}}}$ (quotient d'équivalents : licite).
  \item $n!=\sqrt{2\pi n}\,\Bigl(\dfrac{n}{\mathrm{e}}\Bigr)^{n}(1+\varepsilon_n)$
    avec $\varepsilon_n\to0$, donc
    $\ln(n!)=n\ln n-n+\dfrac12\ln(2\pi n)+\ln(1+\varepsilon_n)$, et le
    dernier terme tend vers $0$ : c'est le $o(1)$ annoncé.
  \item $\dfrac{\ln(n!)}{n}=\ln n-1+\dfrac{\ln(2\pi n)}{2n}+o\Bigl(\dfrac1n\Bigr)
    =\ln\dfrac{n}{\mathrm{e}}+o(1)$, donc
    $\dfrac{(n!)^{1/n}}{n/\mathrm{e}}=\mathrm{e}^{o(1)}\to1$ :
    $\res{(n!)^{1/n}\sim\dfrac{n}{\mathrm{e}}}$.
\end{questions}
\end{corrige}
\end{blocLivre}

\begin{blocLivre}{cor-11-15}
\begin{corrige}{15}{Le développement de $(1+1/n)^{n}$}
\demarche{Développer le logarithme à l'ordre $1/n^{2}$, puis composer par
l'exponentielle.}
\begin{questions}
  \item $\ln u_n=n\ln\Bigl(1+\dfrac1n\Bigr)
    =n\Bigl(\dfrac1n-\dfrac{1}{2n^{2}}+\dfrac{1}{3n^{3}}+o\Bigl(\dfrac{1}{n^{3}}\Bigr)\Bigr)
    =1-\dfrac{1}{2n}+\dfrac{1}{3n^{2}}+o\Bigl(\dfrac{1}{n^{2}}\Bigr)$.
  \item $u_n=\mathrm{e}\cdot\mathrm{e}^{v_n}$ avec
    $v_n=-\dfrac{1}{2n}+\dfrac{1}{3n^{2}}+o\Bigl(\dfrac{1}{n^{2}}\Bigr)$, et
    $\mathrm{e}^{v_n}=1+v_n+\dfrac{v_n^{2}}{2}+o\bigl(v_n^{2}\bigr)$ avec
    $\dfrac{v_n^{2}}{2}=\dfrac{1}{8n^{2}}+o\Bigl(\dfrac{1}{n^{2}}\Bigr)$ :
    comme $\dfrac13+\dfrac18=\dfrac{11}{24}$,
    \[
    u_n=\res{\mathrm{e}\Bigl(1-\frac{1}{2n}+\frac{11}{24n^{2}}
      +o\Bigl(\frac{1}{n^{2}}\Bigr)\Bigr)}.
    \]
    Alors $n(\mathrm{e}-u_n)=\dfrac{\mathrm{e}}{2}-\dfrac{11\mathrm{e}}{24n}
    +o\Bigl(\dfrac1n\Bigr)\to\res{\dfrac{\mathrm{e}}{2}}$ — c'est la version
    « suites » de la limite b) de l'exercice 9.
\end{questions}
\end{corrige}
\end{blocLivre}

\begin{blocLivre}{cor-11-16}
\begin{corrige}{16}{Une réciproque au microscope}
\demarche{Réinjecter l'estimation grossière dans l'équation pour la raffiner :
c'est le principe de tout développement asymptotique.}
\begin{questions}
  \item $f:x\mapsto x+\ln x$ est une bijection strictement croissante de
    $\intOO{0}{+\infty}$ sur $\R$, donc $x_y=f^{-1}(y)$ est croissante ;
    si elle était majorée par $M$, on aurait $y=f(x_y)\leqslant f(M)$ pour
    tout $y$ : absurde. Donc $x_y\to+\infty$. Alors
    $y=x_y\Bigl(1+\dfrac{\ln x_y}{x_y}\Bigr)$ et
    $\dfrac{\ln x_y}{x_y}\to0$ (croissances comparées), d'où
    $\res{x_y\sim y}$.
  \item De $x_y=y-\ln x_y$ et $\dfrac{x_y}{y}\to1$ :
    $\ln x_y=\ln y+\ln\dfrac{x_y}{y}=\ln y+o(1)$, donc
    \[
    \res{x_y=y-\ln y+o(1)}
    \]
    — mieux que le $o(\ln y)$ demandé : réinjecter a transformé
    l'équivalent $x_y\sim y$ en un développement à deux termes.
\end{questions}
\end{corrige}
\end{blocLivre}

\begin{blocLivre}{cor-11-17}
\begin{corrige}{17}{Les solutions de $\tan x=x$}
\demarche{Existence par monotonie, localisation, puis poser l'écart au bord
et le développer.}
\begin{questions}
  \item Sur $J_n=\bigl]n\pi-\frac{\pi}{2},\,n\pi+\frac{\pi}{2}\bigr[$, la
    fonction $g(x)=\tan x-x$ vérifie $g'(x)=\tan^{2}x\geqslant0$, nul au
    seul point $n\pi$ : $g$ est strictement croissante, de limites
    $-\infty$ et $+\infty$ aux bords, d'où un unique zéro $x_n$. Comme
    $g(n\pi)=-n\pi<0$, on a $x_n\in\bigl]n\pi,\,n\pi+\frac{\pi}{2}\bigr[$.
  \item L'encadrement $n\pi<x_n<n\pi+\dfrac{\pi}{2}$ donne
    $\dfrac{x_n}{n\pi}\to1$ : $\res{x_n\sim n\pi}$. Posons
    $\varepsilon_n=n\pi+\dfrac{\pi}{2}-x_n\in\bigl]0,\frac{\pi}{2}\bigr[$ :
    \[
    x_n=\tan x_n=\tan\Bigl(\frac{\pi}{2}-\varepsilon_n\Bigr)
    =\frac{1}{\tan\varepsilon_n},
    \qquad\text{donc}\qquad
    \tan\varepsilon_n=\frac{1}{x_n}\to0 .
    \]
    Ainsi $\varepsilon_n\to0$, puis
    $\varepsilon_n\sim\tan\varepsilon_n=\dfrac{1}{x_n}\sim\dfrac{1}{n\pi}$ :
    \[
    \res{x_n=n\pi+\frac{\pi}{2}-\frac{1}{n\pi}+o\Bigl(\frac1n\Bigr)}.
    \]
\end{questions}
\end{corrige}
\end{blocLivre}

\begin{blocLivre}{cor-11-18}
\begin{corrige}{18}{Une suite d'intégrales}
\demarche{Une relation exacte (b) plus une inégalité de monotonie : l'équivalent
tombe tout seul.}
\begin{questions}
  \item Intégrande positif sur $\intFF{0}{1}$ : $I_n\geqslant0$ ;
    $x^{n+1}\leqslant x^{n}$ sur $\intFF{0}{1}$ : $(I_n)$ décroît ; enfin
    $I_n\leqslant\displaystyle\int_{0}^{1}x^{n}\mathrm{d}x=\dfrac{1}{n+1}\to0$.
  \item $I_n+I_{n+1}=\displaystyle\int_{0}^{1}\frac{x^{n}(1+x)}{1+x}\,\mathrm{d}x
    =\int_{0}^{1}x^{n}\,\mathrm{d}x=\res{\dfrac{1}{n+1}}$. Par décroissance,
    $2I_{n+1}\leqslant I_n+I_{n+1}\leqslant2I_n$, d'où
    $\dfrac{1}{2(n+1)}\leqslant I_n\leqslant\dfrac{1}{2n}$ ($n\geqslant1$),
    et donc $2nI_n\to1$ : $\res{I_n\sim\dfrac{1}{2n}}$.
  \item De $I_k+I_{k-1}=\dfrac1k$, en posant $b_k=(-1)^{k-1}I_k$ :
    $(-1)^{k-1}\bigl(I_k+I_{k-1}\bigr)=b_k-b_{k-1}$, et la somme télescope :
    \[
    \sum_{k=1}^{n}\frac{(-1)^{k-1}}{k}=b_n-b_0=\ln2+(-1)^{n-1}I_n ,
    \]
    puisque $I_0=\displaystyle\int_{0}^{1}\frac{\mathrm{d}x}{1+x}=\ln2$. Comme
    $I_n\to0$, les sommes alternées convergent vers $\res{\ln2}$, avec une
    erreur $I_n\sim\dfrac{1}{2n}$ — lente, mais parfaitement quantifiée.
\end{questions}
\end{corrige}
\end{blocLivre}

\begin{blocLivre}{cor-11-19}
\begin{corrige}{19}{Moyennes généralisées}
\demarche{Tout passe par le logarithme, développé à l'ordre 1 en $t$.}
\begin{questions}
  \item $m(1)=\dfrac{a+b}{2}$ : moyenne arithmétique ;
    $m(-1)=\Bigl(\dfrac{a^{-1}+b^{-1}}{2}\Bigr)^{-1}=\dfrac{2ab}{a+b}$ :
    moyenne harmonique.
  \item $a^{t}=\mathrm{e}^{t\ln a}=1+t\ln a+o(t)$, donc
    $\dfrac{a^{t}+b^{t}}{2}=1+t\,\dfrac{\ln a+\ln b}{2}+o(t)$ et
    \[
    \ln m(t)=\frac1t\ln\Bigl(1+t\ln\sqrt{ab}+o(t)\Bigr)
    =\ln\sqrt{ab}+o(1)\xrightarrow[t\to0]{}\ln\sqrt{ab} :
    \]
    $m(t)\to\res{\sqrt{ab}}$, la moyenne géométrique s'intercale donc entre
    harmonique ($t=-1$) et arithmétique ($t=1$).
  \item Pour $t>0$ : $m(t)=b\Bigl(\dfrac{1+(a/b)^{t}}{2}\Bigr)^{1/t}$ avec
    $0<\dfrac{a}{b}<1$, donc $\Bigl(\dfrac{a}{b}\Bigr)^{t}\to0$ et
    $\dfrac1t\ln\dfrac{1+(a/b)^{t}}{2}\to0$ (numérateur borné) :
    $m(t)\to\res{b=\max(a,b)}$. En $-\infty$, le même calcul avec $a$ en
    facteur donne $m(t)\to\res{a=\min(a,b)}$.
\end{questions}
\end{corrige}
\end{blocLivre}

\begin{blocLivre}{cor-11-20}
\begin{corrige}{20}{Héron : la convergence quadratique quantifiée}
\demarche{L'erreur se recopie au carré : une récurrence sur l'exposant $2^{n}$.}
\begin{questions}
  \item $u_{n+1}-\sqrt2=\dfrac{u_n^{2}-2\sqrt2\,u_n+2}{2u_n}
    =\dfrac{(u_n-\sqrt2)^{2}}{2u_n}=\dfrac{e_n^{2}}{2u_n}
    \leqslant\dfrac{e_n^{2}}{2\sqrt2}$, car $u_n\geqslant\sqrt2$.
  \item Initialisation : $e_0=2\sqrt2\,\rho$ par définition de $\rho$.
    Hérédité : si $e_n\leqslant2\sqrt2\,\rho^{2^{n}}$, alors
    $e_{n+1}\leqslant\dfrac{\bigl(2\sqrt2\,\rho^{2^{n}}\bigr)^{2}}{2\sqrt2}
    =2\sqrt2\,\rho^{2^{n+1}}$ : c'est la formule au rang $n+1$.
  \item Ici $\rho=\dfrac{2-\sqrt2}{2\sqrt2}=\dfrac{\sqrt2-1}{2}\approx0{,}207$,
    donc $-\log_{10}e_n\geqslant2^{n}\times0{,}68-0{,}46$ : le nombre de
    décimales exactes \emph{double} au moins à chaque itération. Pour $n=4$ :
    $e_4\leqslant2\sqrt2\,\rho^{16}\approx\res{3\cdot10^{-11}}$ — quatre
    itérations suffisent pour la précision d'une calculatrice.
\end{questions}
\end{corrige}
\end{blocLivre}

\begin{blocLivre}{cor-11-21}
\begin{corrige}{21}{Le grand classique $u_{n+1}=\sin u_n$}
\demarche{Quand $u_{n+1}-u_n$ est trop fin, on cherche la bonne puissance :
ici $1/u_n^{2}$ progresse arithmétiquement, et Cesàro conclut.}
\begin{questions}
  \item Sur $\bigl]0,\frac{\pi}{2}\bigr]$, $0<\sin u<u$ : par récurrence
    $u_n\in\intOO{0}{1}$ pour $n\geqslant1$ et $(u_n)$ décroît. Minorée par
    $0$, elle converge vers $\ell$ avec $\sin\ell=\ell$, donc $\ell=0$.
  \item $\sin u=u-\dfrac{u^{3}}{6}+o\bigl(u^{3}\bigr)$ donne
    $\sin^{2}u=u^{2}-\dfrac{u^{4}}{3}+o\bigl(u^{4}\bigr)$, d'où
    \[
    \frac{1}{u_{n+1}^{2}}-\frac{1}{u_n^{2}}
    =\frac{u_n^{2}-\sin^{2}u_n}{u_n^{2}\sin^{2}u_n}
    =\frac{\frac{u_n^{4}}{3}+o\bigl(u_n^{4}\bigr)}{u_n^{4}+o\bigl(u_n^{4}\bigr)}
    \xrightarrow[n\to+\infty]{}\frac13 .
    \]
  \item Par Cesàro (chapitre 5), la moyenne des différences tend vers la
    même limite :
    $\dfrac{1}{n}\Bigl(\dfrac{1}{u_n^{2}}-\dfrac{1}{u_0^{2}}\Bigr)\to\dfrac13$,
    donc $\dfrac{1}{u_n^{2}}\sim\dfrac{n}{3}$, puis $u_n^{2}\sim\dfrac3n$ et,
    par racine carrée (puissance fixe),
    $\res{u_n\sim\sqrt{\dfrac3n}}$.
\end{questions}
\end{corrige}
\end{blocLivre}

\begin{blocLivre}{cor-11-22}
\begin{corrige}{22}{Étude asymptotique complète}
\demarche{Limites, variations, développement en $1/x$ : les trois étages
d'une étude de fonction asymptotique.}
\begin{questions}
  \item Avec $u=\dfrac1x$ : en $0^{+}$, $f(x)=\dfrac{\mathrm{e}^{u}}{u}\to+\infty$
    (croissances comparées) ; en $0^{-}$, $u\to-\infty$ et
    $\dfrac{\mathrm{e}^{u}}{u}\to0$ : $f$ se prolonge à gauche par $f(0)=0$,
    et $\dfrac{f(x)-0}{x}=\mathrm{e}^{1/x}\to0$ : dérivée à gauche nulle —
    la courbe arrive \emph{tangentiellement} à l'axe. En $\pm\infty$,
    $f(x)\sim x\to\pm\infty$.
  \item $f'(x)=\mathrm{e}^{1/x}\Bigl(1-\dfrac1x\Bigr)
    =\mathrm{e}^{1/x}\,\dfrac{x-1}{x}$ : sur $\intOO{0}{+\infty}$, $f$
    décroît sur $\intOF{0}{1}$ puis croît, avec un minimum
    $f(1)=\mathrm{e}$ ; sur $\intOO{-\infty}{0}$, $\dfrac{x-1}{x}>0$ : $f$
    est croissante.
  \item $f(x)=x\Bigl(1+\dfrac1x+\dfrac{1}{2x^{2}}+o\Bigl(\dfrac{1}{x^{2}}\Bigr)\Bigr)
    =\res{x+1+\dfrac{1}{2x}+o\Bigl(\dfrac1x\Bigr)}$ : la droite $y=x+1$ est
    asymptote des deux côtés, courbe au-dessus en $+\infty$
    ($\tfrac{1}{2x}>0$), au-dessous en $-\infty$.
\end{questions}
\medskip
\begin{center}
\begin{tikzpicture}[scale=0.52,line cap=round]
  \clip (-3.9,-3.1) rectangle (3.9,5.9);
  \draw[otMuted,line width=0.5pt,-{Stealth[length=1.6mm]}] (-3.7,0) -- (3.7,0);
  \draw[otMuted,line width=0.5pt,-{Stealth[length=1.6mm]}] (0,-2.9) -- (0,5.6);
  \draw[otGreenDark,line width=0.9pt,dash pattern=on 3pt off 1.8pt]
    plot[domain=-3.7:3.7,samples=30] ({\x},{\x+1});
  \draw[otNavy,line width=1.3pt]
    plot[domain=0.42:3.7,samples=120] ({\x},{\x*exp(1/\x)});
  \draw[otNavy,line width=1.3pt]
    plot[domain=-3.7:-0.12,samples=120] ({\x},{\x*exp(1/\x)});
  \fill[otGreenDark] (1,2.71828) circle (1.6pt);
  \node[font=\tiny,text=otGreenDark,anchor=west] at (1.1,2.55) {$(1,\mathrm{e})$};
  \node[font=\tiny,text=otGreenDark,rotate=38] at (-2.15,-1.65) {$y=x+1$};
  \node[font=\tiny,text=otNavy,anchor=west] at (2.15,4.95) {$y=x\,\mathrm{e}^{1/x}$};
\end{tikzpicture}\\[0.5mm]
\legende{La courbe de $x\,\mathrm{e}^{1/x}$ : minimum $(1,\mathrm{e})$,
asymptote $y=x+1$ des deux côtés, arrivée tangente à l'axe en $0^{-}$.}
\end{center}
\end{corrige}
\end{blocLivre}

\begin{blocLivre}{cor-11-23}
\begin{corrige}{23}{Wallis, sans Stirling}
\demarche{Un invariant exact, un encadrement de monotonie : l'équivalent de
$I_n$ n'exige aucun théorème admis.}
\begin{questions}
  \item Pour $n\geqslant2$, la relation $I_n=\dfrac{n-1}{n}I_{n-2}$ donne
    $n\,I_n\,I_{n-1}=(n-1)\,I_{n-1}\,I_{n-2}$ : la suite est constante,
    égale à $1\cdot I_1\cdot I_0=1\cdot1\cdot\dfrac{\pi}{2}
    =\res{\dfrac{\pi}{2}}$.
  \item Par décroissance, $I_{n+1}\leqslant I_n\leqslant I_{n-1}$, donc
    $\dfrac{n}{n+1}=\dfrac{I_{n+1}}{I_{n-1}}\leqslant\dfrac{I_{n+1}}{I_n}\leqslant1$ :
    par encadrement, $\dfrac{I_{n+1}}{I_n}\to1$. Alors
    $\dfrac{\pi}{2}=n\,I_n\,I_{n-1}\sim n\,I_n^{2}$, et $I_n>0$ :
    $\res{I_n\sim\sqrt{\dfrac{\pi}{2n}}}$.
  \item La récurrence paire donne
    $I_{2n}=\dfrac{2n-1}{2n}\cdot\dfrac{2n-3}{2n-2}\cdots\dfrac12\cdot I_0
    =\dfrac{(2n)!}{4^{n}(n!)^{2}}\cdot\dfrac{\pi}{2}
    =\dfrac{\pi}{2}\cdot\dfrac{1}{4^{n}}\binom{2n}{n}$. En combinant avec
    b) : $\dfrac{1}{4^{n}}\binom{2n}{n}=\dfrac{2}{\pi}I_{2n}
    \sim\dfrac{2}{\pi}\sqrt{\dfrac{\pi}{4n}}=\dfrac{1}{\sqrt{\pi n}}$, soit
    \[
    \res{\binom{2n}{n}\sim\frac{4^{n}}{\sqrt{\pi n}}}
    \]
    — le résultat de l'exercice 14 a), obtenu cette fois sans admettre
    Stirling : Wallis en contient déjà la constante $\sqrt{\pi}$.
\end{questions}
\end{corrige}
\end{blocLivre}

\begin{blocLivre}{cor-11-24}
\begin{corrige}{24}{Vers Stirling : encadrer $\ln(n!)$}
\demarche{Comparer la somme $\sum\ln k$ à l'intégrale de $\ln$ : la méthode
série-intégrale dans sa forme la plus simple.}
\begin{questions}
  \item $\ln$ est croissante : pour $t\in\intFF{k}{k+1}$,
    $\ln k\leqslant\ln t\leqslant\ln(k+1)$ ; en intégrant sur un intervalle
    de longueur $1$, l'encadrement annoncé.
  \item En sommant de $k=1$ à $n-1$, avec
    $\displaystyle\int_{1}^{n}\ln t\,\mathrm{d}t=n\ln n-n+1$ :
    \[
    \ln\bigl((n-1)!\bigr)\leqslant n\ln n-n+1\leqslant\ln(n!) .
    \]
    La seconde inégalité est la minoration voulue ; pour la majoration,
    $\ln(n!)=\ln\bigl((n-1)!\bigr)+\ln n\leqslant n\ln n-n+1+\ln n$.
  \item En divisant par $n$ :
    $\ln n-1+\dfrac1n\leqslant\dfrac{\ln(n!)}{n}\leqslant\ln n-1+\dfrac{1+\ln n}{n}$,
    donc $\dfrac{\ln(n!)}{n}-\ln\dfrac{n}{\mathrm{e}}\to0$ et
    $\dfrac{(n!)^{1/n}}{n/\mathrm{e}}=\mathrm{e}^{o(1)}\to1$ :
    $\res{(n!)^{1/n}\sim\dfrac{n}{\mathrm{e}}}$. Stirling précise ce que
    notre fenêtre de largeur $\ln n$ laisse flou : le terme manquant vaut
    exactement $\dfrac12\ln(2\pi n)+o(1)$.
\end{questions}
\end{corrige}
\end{blocLivre}

% !TEX root = ../../main.tex
%% Corrigé du défi du chapitre 11 — vérifié par verifications/defis_11_15.py.

\begin{blocLivre}{cor-11-25a}
\begin{corrige}{25}{La fonction W de Lambert}
\demarche{La réciproque d'une bijection se dérive par la formule
$1/g'(W)$ ; son DL s'obtient par identification dans l'équation qui la
définit ; son comportement à l'infini, en réinjectant l'estimation dans
$W+\ln W=\ln y$.}
\partieExo{A}{Définition et dérivée}
\begin{questions}
  \item $g'(x)=(1+x)\,\mathrm{e}^{x}>0$ sur $\intOO{-1}{+\infty}$ : $g$ est
    continue, strictement croissante, $g(-1)=-\frac1{\mathrm{e}}$ et
    $g(x)\to+\infty$. Par le théorème de la bijection,
    $g:\intFO{-1}{+\infty}\to\intFO{-\frac1{\mathrm{e}}}{+\infty}$ est
    bijective. $g(0)=0$, $g(1)=\mathrm{e}$, $g(-1)=-\frac1{\mathrm{e}}$
    donnent $\res{W(0)=0}$, $\res{W(\mathrm{e})=1}$,
    $\res{W\bigl(-\tfrac1{\mathrm{e}}\bigr)=-1}$.
  \item $W$ est continue et strictement croissante (réciproque d'une
    bijection continue strictement croissante). Comme $g'$ ne s'annule pas
    sur $\intOO{-1}{+\infty}$, $W$ est dérivable sur
    $\intOO{-\frac1{\mathrm{e}}}{+\infty}$ et
    $W'=\dfrac{1}{g'(W)}=\dfrac{1}{(1+W)\mathrm{e}^{W}}$. Pour $y\neq0$,
    $\mathrm{e}^{W}=\frac yW$ d'où la seconde forme. $\res{W'(0)=1}$. En
    $-\frac1{\mathrm{e}}$, $g'(-1)=0$ : $W'(y)\to+\infty$, \emph{tangente
    verticale}.
\end{questions}
\partieExo{B}{Développement limité en $0$}
\begin{questions}[resume]
  \item Écrivons $W=y+a_2y^{2}+a_3y^{3}+o\bigl(y^{3}\bigr)$ ($W(0)=0$,
    $W'(0)=1$). Comme $W=O(y)$,
    $W\mathrm{e}^{W}=W+W^{2}+\frac{W^{3}}{2}+o\bigl(y^{3}\bigr)$, avec
    $W^{2}=y^{2}+2a_2y^{3}+o\bigl(y^{3}\bigr)$ et $W^{3}=y^{3}+o\bigl(y^{3}\bigr)$ :
    \[
    W\mathrm{e}^{W}=y+(a_2+1)y^{2}+\Bigl(a_3+2a_2+\frac12\Bigr)y^{3}+o\bigl(y^{3}\bigr)=y .
    \]
    D'où $a_2=-1$, $a_3=\frac32$ :
    $\res{W(y)=y-y^{2}+\tfrac32y^{3}+o\bigl(y^{3}\bigr)}$.
  \item À l'ordre $4$, avec $W=y-y^{2}+\frac32y^{3}+a_4y^{4}$ :
    $W^{2}=y^{2}-2y^{3}+4y^{4}$, $W^{3}=y^{3}-3y^{4}$, $W^{4}=y^{4}$
    (modulo $o(y^{4})$), et le coefficient de $y^{4}$ dans
    $W+W^{2}+\frac{W^{3}}2+\frac{W^{4}}6$ est
    $a_4+4-\frac32+\frac16=0$, soit $\res{a_4=-\tfrac83}$. Contrôle :
    $\frac{(-n)^{n-1}}{n!}$ vaut $1,-1,\frac32,-\frac83$ pour
    $n=1,2,3,4$.
  \item $x=W(0{,}01)\approx0{,}01-0{,}0001+0{,}0000015=\res{0{,}0099015}$ ;
    le terme suivant, $-\frac83\cdot10^{-8}$, donne l'ordre de grandeur de
    l'erreur ($3\cdot10^{-8}$).
\end{questions}
\end{corrige}
\end{blocLivre}

\begin{blocLivre}{cor-11-25b}
\begin{corrigeSuite}
\partieExo{C}{Comportement en $+\infty$}
\begin{questions}[start=6]
  \item $W$ est croissante ; si elle était majorée par $M$, on aurait
    $y=g(W(y))\leqslant g(M)$ pour tout $y$ : absurde. Donc
    $W(y)\to+\infty$. Pour $y>0$, $W(y)>0$ et $\ln$ de $W\mathrm{e}^{W}=y$
    donne $\res{W+\ln W=\ln y}$. Comme $\ln W=o(W)$ :
    $W(1+o(1))=\ln y$, soit $\res{W(y)\sim\ln y}$.
  \item $\ln W=\ln\bigl(\ln y\,(1+o(1))\bigr)=\ln\ln y+o(1)$, donc
    $W=\ln y-\ln W=\ln y-\ln\ln y+o(1)$. Réinjectons :
    $W=\ln y\Bigl(1-\frac{\ln\ln y}{\ln y}+o\bigl(\tfrac{1}{\ln y}\bigr)\Bigr)$,
    donc
    \[
    \ln W=\ln\ln y-\frac{\ln\ln y}{\ln y}+o\Bigl(\frac{\ln\ln y}{\ln y}\Bigr),
    \]
    et $W=\ln y-\ln W$ donne
    $\res{W(y)=\ln y-\ln\ln y+\frac{\ln\ln y}{\ln y}+o\Bigl(\frac{\ln\ln y}{\ln y}\Bigr)}$.
  \item $x^{x}=y\iff x\ln x=\ln y$. La fonction $x\mapsto x\ln x$ est
    strictement croissante de $\intOO{1}{+\infty}$ sur $\intOO{0}{+\infty}$ :
    unique solution $x_y>1$. Avec $u=\ln x>0$ : $u\,\mathrm{e}^{u}=\ln y$,
    donc $u=W(\ln y)$ et
    $x_y=\mathrm{e}^{W(\ln y)}=\dfrac{\ln y}{W(\ln y)}$ (car
    $\mathrm{e}^{W(t)}=\frac{t}{W(t)}$). Avec la question 6 :
    $\res{x_y\sim\dfrac{\ln y}{\ln\ln y}}$. Pour $y=10^{20}$ :
    $W(\ln y)\approx2{,}800$ et $x_y\approx16{,}45$, alors que
    l'équivalent donne seulement $12{,}0$ : la convergence est très lente,
    d'où l'intérêt du développement plus fin de la question 7.
\end{questions}
\end{corrigeSuite}
\end{blocLivre}

